\chapter{Foundations of Generative Modeling}

\section{Density Estimation and Sampling}

Generative modeling aims to learn the underlying distribution of data and generate new samples from it.

\medskip

\noindent
\textbf{Objective.}  
Given data $x \sim p_{\text{data}}$, learn a model $p_\theta(x)$ that approximates the true distribution.

\medskip

\noindent
\textbf{Density estimation.}
\begin{itemize}
    \item Explicit models define a probability density $p_\theta(x)$
    \item Learning is typically based on maximum likelihood:
\end{itemize}

\[
\max_\theta \; \mathbb{E}_{x \sim p_{\text{data}}}[\log p_\theta(x)].
\]

\medskip

\noindent
\textbf{Sampling.}  
Once a model is learned, we generate new samples:
\[
x \sim p_\theta(x).
\]

\medskip

\noindent
\textbf{Two perspectives.}
\begin{itemize}
    \item \textbf{Likelihood-based:} directly model $p_\theta(x)$
    \item \textbf{Implicit:} define a sampling procedure without explicit density
\end{itemize}

\medskip

\noindent
\textbf{Challenges.}
\begin{itemize}
    \item High-dimensional data
    \item Complex distributions
    \item Efficient sampling
\end{itemize}

\medskip

\noindent
\textbf{Insight.}  
Generative modeling involves both representing distributions and generating samples from them.

\section{Latent Variable Models}

Many generative models introduce latent variables to capture underlying structure.

\medskip

\noindent
\textbf{Latent variable formulation.}
\[
p_\theta(x) = \int p_\theta(x \mid z) \, p(z) \, dz,
\]
where $z$ is a latent variable.

\medskip

\noindent
\textbf{Interpretation.}
\begin{itemize}
    \item $z$ represents hidden factors of variation
    \item $p(z)$ is a simple prior (e.g., Gaussian)
\end{itemize}

\medskip

\noindent
\textbf{Generative process.}
\begin{enumerate}
    \item Sample $z \sim p(z)$
    \item Generate $x \sim p_\theta(x \mid z)$
\end{enumerate}

\medskip

\noindent
\textbf{Learning challenge.}  
The marginal likelihood involves an intractable integral.

\medskip

\noindent
\textbf{Inference.}  
We introduce an approximate posterior $q_\phi(z \mid x)$.

\medskip

\noindent
\textbf{Insight.}  
Latent variables provide a structured way to model complex data distributions.

\section{Energy-Based Perspectives}

An alternative view models probability distributions through energy functions.

\medskip

\noindent
\textbf{Energy-based model (EBM).}
\[
p_\theta(x) = \frac{1}{Z(\theta)} \exp(-E_\theta(x)),
\]
where $E_\theta(x)$ is an energy function and $Z(\theta)$ is the partition function:
\[
Z(\theta) = \int \exp(-E_\theta(x)) \, dx.
\]

\medskip

\noindent
\textbf{Interpretation.}
\begin{itemize}
    \item Low energy $\rightarrow$ high probability
    \item High energy $\rightarrow$ low probability
\end{itemize}

\medskip

\noindent
\textbf{Learning.}
\begin{itemize}
    \item Requires estimating gradients involving $Z(\theta)$
    \item Often uses sampling-based methods
\end{itemize}

\medskip

\noindent
\textbf{Connection to physics.}  
Energy-based models resemble distributions in statistical mechanics.

\medskip

\noindent
\textbf{Insight.}  
Energy-based formulations provide a flexible way to define distributions without requiring normalized densities in closed form.

\section{Evaluation of Generative Models}

Evaluating generative models is challenging due to the complexity of high-dimensional distributions.

\medskip

\noindent
\textbf{Likelihood-based evaluation.}
\begin{itemize}
    \item Log-likelihood
    \item Perplexity (for sequences)
\end{itemize}

\medskip

\noindent
\textbf{Sample quality.}
\begin{itemize}
    \item Visual inspection (e.g., images)
    \item Human evaluation (e.g., text)
\end{itemize}

\medskip

\noindent
\textbf{Divergence measures.}
\begin{itemize}
    \item Kullback--Leibler divergence
    \item Jensen--Shannon divergence
\end{itemize}

\medskip

\noindent
\textbf{Distributional metrics.}
\begin{itemize}
    \item Inception Score
    \item Fr\'echet Inception Distance (FID)
\end{itemize}

\medskip

\noindent
\textbf{Challenges.}
\begin{itemize}
    \item Tradeoff between sample quality and diversity
    \item Difficulty of estimating likelihood in high dimensions
\end{itemize}

\medskip

\noindent
\textbf{Mode collapse.}  
Some models generate limited diversity, focusing on a subset of possible outputs.

\medskip

\noindent
\textbf{Insight.}  
No single metric fully captures the quality of a generative model.

\section{A Unifying Perspective}

Generative modeling can be understood through multiple complementary frameworks:

\begin{itemize}
    \item \textbf{Probabilistic:} modeling $p(x)$ directly
    \item \textbf{Latent variable:} introducing hidden structure
    \item \textbf{Energy-based:} defining distributions via energy
\end{itemize}

\medskip

\noindent
\textbf{Core challenges.}
\begin{itemize}
    \item High-dimensional modeling
    \item Efficient sampling
    \item Reliable evaluation
\end{itemize}

\medskip

\noindent
\textbf{Key insight.}  
Generative models aim to capture the structure of data distributions in a way that enables both understanding and synthesis.

\section{Outlook}

This chapter introduced the foundations of generative modeling:
\begin{itemize}
    \item density estimation and sampling,
    \item latent variable models,
    \item energy-based perspectives,
    \item evaluation of generative models.
\end{itemize}

\medskip

In the next chapters, we study specific classes of generative models, including variational autoencoders, generative adversarial networks, and diffusion models.

\medskip

\noindent
\textbf{Guiding principle.}  
Generative modeling seeks not only to predict data, but to understand and reproduce its underlying structure.